% META: {"id":"1SPE-PROBCOND-EX-029","chapitre":"1SPE-PROBA-COND","type_objet":"exercice","capacites":["1SPE-PROBA-COND-C3"],"capacites_codes":["C3"],"methodes":["M3"],"parcours":3,"competences":["calculer","raisonner"],"duree_min":15,"mode_creation":"ex_nihilo","sources_inspiration":[],"parametres_sympy":{},"fichier_tex":"chapitres/1SPE-PROBA-COND/exercices/1SPE-PROBCOND-EX-029.tex","corrige_tex":"chapitres/1SPE-PROBA-COND/corriges/1SPE-PROBCOND-CO-029.tex","status":"approved","origin":"generated","status_history":["generated","audited","accepted","approved"],"release_acceptance":"RELEASE_OWNER_BATCH_ACCEPTANCE","acceptance_closure_digest":"sha256:9b3ccf9a81c5520fbb7e03b7b2d3e7bf2057a02834b6d908bf3be5f49f3b553f","acceptance_receipt_digest":"sha256:4fad86f0282e0fa4e0419cca1ad6379ae7af5a280c04a4a90880f90a1c044242"}
% BEGIN-VERIFY
% from sympy import Rational, Symbol, solve
% p = Symbol('p')
% # P(B)=p, P_B(A)=1/2, P_Bbar(A)=1/5, P(A)=3/10
% eq = p * Rational(1,2) + (1-p) * Rational(1,5) - Rational(3,10)
% sol = solve(eq, p)
% assert sol[0] == Rational(1, 3)
% END-VERIFY
\begin{exercice}{1SPE-PROBCOND-EX-029}{3}{15}
On donne $P_B(A) = \dfrac{1}{2}$, $P_{\overline{B}}(A) = \dfrac{1}{5}$ et $P(A) = \dfrac{3}{10}$.

Déterminer $P(B)$.
\end{exercice}
